\documentclass[12pt,a4paper]{article}
\usepackage[margin=25mm, headheight=15pt, headsep=10pt]{geometry}
\usepackage{fontspec}
\usepackage[table]{xcolor} % Note: table option is required for rowcolors
\usepackage{titlesec}
\usepackage{tocloft}
\usepackage{fancyhdr}
\usepackage{booktabs}
\usepackage{tcolorbox}
\tcbuselibrary{skins, breakable}
\usepackage[colorlinks=true, linkcolor=emerald, urlcolor=emerald, bookmarks=true]{hyperref}

\definecolor{emerald}{RGB}{0,102,51}
\definecolor{darkred}{RGB}{139,0,0}
\definecolor{lightgray}{RGB}{245,245,245}

\usepackage[nil]{babel}
\babelprovide[import, main]{bengali}
\babelprovide[import]{arabic}
\babelfont{rm}[Renderer=HarfBuzz,Scale=1.0]{Noto Serif Bengali}
\babelfont{sf}[Renderer=HarfBuzz]{Noto Sans Bengali}
\babelfont{tt}[Renderer=HarfBuzz]{Noto Sans Bengali}
\babelfont[arabic]{rm}[Renderer=HarfBuzz,Scale=1.1]{Noto Naskh Arabic}

% Section styling
\titleformat{\section}{\Large\bfseries\color{emerald}}{\thesection.}{0.5em}{}
\titleformat{\subsection}{\large\bfseries\color{emerald!85!black}}{\thesubsection.}{0.5em}{}
\titleformat{\subsubsection}{\normalsize\bfseries\color{emerald!70!black}}{\thesubsubsection.}{0.5em}{}
\titlespacing*{\section}{0pt}{18pt}{8pt}

% Fancy Header/Footer setup
\pagestyle{fancy}
\fancyhf{} % clear all header and footer fields
\fancyhead[L]{\color{emerald}\small\textbf{\nouppercase{\leftmark}}}
\fancyfoot[L]{\scriptsize\color{black!50}{\lat{AI}}-সংকলিত, আলেম-যাচাইকৃত নয়}
\fancyfoot[C]{\color{emerald!70!black}\thepage}
\renewcommand{\headrulewidth}{0.4pt}
\renewcommand{\headrule}{\hbox to\headwidth{\color{emerald}\leaders\hrule height \headrulewidth\hfill}}
\renewcommand{\footrulewidth}{0pt}

% Custom boxes
% 1. Ayat/Hadith Box
\newtcolorbox{ayatbox}[1][]{
    enhanced, breakable, 
    colback=emerald!5, 
    colframe=emerald, 
    boxrule=0.5pt, 
    arc=3pt, 
    left=10pt, right=10pt, top=10pt, bottom=10pt,
    #1
}

% 2. Quote/Translation Box
\newtcolorbox{quotebox}[1][]{
    enhanced, breakable,
    frame hidden,
    colback=lightgray,
    borderline west={3pt}{0pt}{emerald},
    left=15pt, right=10pt, top=10pt, bottom=10pt,
    sharp corners,
    #1
}

% 3. AI-generated notice box
\newtcolorbox{warnbox}[1][]{
    enhanced, breakable,
    colback=red!4,
    colframe=darkred,
    boxrule=0.8pt,
    arc=3pt,
    left=10pt, right=10pt, top=10pt, bottom=10pt,
    #1
}

% Commands
\newcommand{\arab}[1]{\foreignlanguage{arabic}{#1}}
\newcommand{\kib}[1]{\textcolor{emerald}{\textbf{#1}}}

% Latin (English) fallback: Noto Serif Bengali has NO Latin glyphs
\newfontfamily\latfont[Renderer=HarfBuzz]{Noto Serif}
\newcommand{\lat}[1]{{\latfont #1}}

\setlength{\parskip}{5pt}
\setlength{\parindent}{0pt}

% Fix for missing bullet in Noto Serif Bengali
\renewcommand{\labelitemi}{$\bullet$}



\begin{document}
\begin{center}{\Large\bfseries\color{emerald} আল্লাহর রহমত ও তাওবা — \lat{04}-আবদুল্লাহ-হিমার-মদপান-ও-নবীজির-নিষেধ}\end{center}
\vspace{1em}
\section*{ঘটনা ৪ — আবদুল্লাহ \lat{“}হিমার\lat{”}: বারবার মদপান, তবু \lat{“}অভিশাপ দিও না\lat{”}}

\begin{warnbox}
 \textbf{গুরুত্বপূর্ণ নোটিশ (\lat{Important} \lat{Notice}):} এই কন্টেন্টটি \lat{AI} (কৃত্রিম বুদ্ধিমত্তা) দিয়ে প্রস্তুত ও সংকলিত। \lat{AI} কঠোর সূত্র-মান নীতি মেনে শুধু নির্ভরযোগ্য উৎস থেকে তথ্য সংগ্রহ করেছে — কেবল সহীহ/হাসান হাদিস ও সহীহ সনদের আসার রাখা হয়েছে, দুর্বল সনদ বাদ দেওয়া হয়েছে। তবে এটি কোনো আলেম বা মুহাদ্দিস কর্তৃক যাচাইকৃত নয়।
\end{warnbox}

\begin{quote}
\textbf{সম্পর্কিত ফাইল:} হাদিস ফাইল — পাপীকে অভিশাপ নয়; ব্যবহারিক গাইড — পাপীর প্রতি দৃষ্টি।
\end{quote}

\medskip\hrule\medskip

\subsection*{১. সূত্র কার্ড (\lat{Source} \lat{Card})}

\begin{table}[h]\centering\small
\begin{tabular}{ll}
বিষয় & বিবরণ \\
\hline
\textbf{ঘটনা} & আবদুল্লাহ (রাঃ) — ডাকনাম \lat{“}হিমার\lat{”} — বারবার মদপানের অপরাধে দণ্ড পেয়েছেন; এক সাহাবি তাকে অভিশাপ করতে চাইলে নবীজি (সা.) নিষেধ করে বললেন — সে আল্লাহ ও তাঁর রাসূলকে ভালোবাসে \\
\textbf{বর্ণনাকারী} & উমার ইবনে খাত্তাব (রাঃ) \\
\textbf{গ্রন্থ ও নম্বর} & সহীহ বুখারি ৬৭৮০; সহীহ মুসলিম ২৭৬৮ \\
\textbf{অধ্যায়} & কিতাবুল হুদুদ — \lat{“}মদ্যপায়ীকে অভিশাপ দেওয়া অপছন্দ — সে কাফির নয়\lat{”} \\
\textbf{সনদের মান} & \textbf{সহীহ} (বুখারি-মুসলিম) \\
\textbf{স্থান ও সময়} & মদিনা — নবীজির (সা.) জীবদ্দশায় \\
\hline
\end{tabular}
\end{table}

\medskip\hrule\medskip

\subsection*{২. সংক্ষিপ্ত পটভূমি (\lat{Background})}

নবীজির (সা.) যুগে মদ হারাম হওয়ার পরও কিছু সাহাবি পুরনো অভ্যাসে বারবার পড়তেন — শাস্তি (দণ্ড) পেতেন, আবার পাপ করতেন। আবদুল্লাহ (রাঃ) এমনই ছিলেন — নবীজিকে (সা.) হাসাতেন, ভালোবাসতেন; কিন্তু মদের সাথে তার লড়াই চলত। এই ঘটনাটি ইসলামের সবচেয়ে গুরুত্বপূর্ণ শিক্ষা দেয় — \textbf{পাপী মুমিনকে হারিয়ে গেছে ভাবা যাবে না; পাপ ও ঈমান একই অন্তরে থাকতে পারে — লড়াই চলতে থাকে।}

\medskip\hrule\medskip

\subsection*{৩. পূর্ণ ঘটনার বিশুদ্ধ বর্ণনা (\lat{The} \lat{Full} \lat{Narration})}

\begin{quote}
নবীজির (সা.) যুগে আবদুল্লাহ নামে এক ব্যক্তি ছিল, যার ডাকনাম ছিল \textbf{\lat{“}হিমার\lat{”}} (গাধা) — সে নবীজিকে (সা.) হাসাতেন। নবীজি (সা.) মদপানের কারণে তাকে দণ্ড দিয়েছেন। একদিন তাকে আবার একই অপরাধে আনা হলো — নবীজি (সা.) আদেশ দিলেন, তাকে দণ্ড দেওয়া হলো। তখন জনতার মধ্যে এক ব্যক্তি বলল: \textbf{\lat{“}হে আল্লাহ! তাকে অভিশাপ দাও — কতবারই না তাকে (এই অপরাধে) আনা হলো!\lat{”}} নবীজি (সা.) বললেন: \textbf{\lat{“}তাকে অভিশাপ দিও না! আল্লাহর কসম! আমি জানি — সে আল্লাহ ও তাঁর রাসূলকে ভালোবাসে।\lat{”}}
\end{quote}

\medskip\hrule\medskip

\subsection*{৪. শব্দের অর্থ (\lat{Key} \lat{Words})}

\begin{table}[h]\centering\small
\begin{tabular}{ll}
শব্দ & অর্থ \\
\hline
\textbf{হিমার (\arab{حمار})} & গাধা — ডাকনাম; এখানে উপহাস নয়, সাধারণ নামকরণ \\
\textbf{জালাদাহু ফিশ শারাব (\arab{جلده} \arab{في} \arab{الشراب})} & মদের অপরাধে তাকে দণ্ড দিয়েছেন \\
\textbf{আলা-নহু (\arab{العنوه})} & তাকে অভিশাপ দাও \\
\textbf{ইউহিব্বুল্লাহা ওয়া রাসূলাহু} & সে আল্লাহ ও তাঁর রাসূলকে ভালোবাসে \\
\hline
\end{tabular}
\end{table}

\medskip\hrule\medskip

\subsection*{৫. সংশ্লিষ্ট দলিল ও রেফারেন্স}

\begin{quote}
\textbf{\lat{“}তাকে অভিশাপ দিও না! আল্লাহর কসম! আমি জানি — সে আল্লাহ ও তাঁর রাসূলকে ভালোবাসে।\lat{”}} — সহীহ বুখারি ৬৭৮০
\end{quote}

\begin{quote}
\textbf{\lat{“}প্রতিটি আদম সন্তান গুনাহ করে, আর গুনাহকারীদের মধ্যে সেরা তারা, যারা তাওবা করে।\lat{”}} — সুনানে ইবনে মাজাহ ৪২৫১ (হাসান)
\end{quote}

\begin{quote}
\textbf{\lat{“}তোমরা রাতে ও দিনে পাপ করতে থাকো, আর আমি সব পাপ ক্ষমা করি — সুতরাং ক্ষমা চাও।\lat{”}} — সহীহ মুসলিম ২৫৭৭ (হাদিস কুদসি)
\end{quote}

\subsubsection*{মূল শিক্ষা (দলিলভিত্তিক)}

\begin{enumerate}
\item \textbf{পাপীকে হারিয়ে গেছে ভাবা নবীজির (সা.) শিক্ষার বিপরীত:} বারবার মদপান — বড় পাপ — তবু নবীজি (সা.) বললেন সে আল্লাহ ও রাসূলকে ভালোবাসে; পাপ ও ভালোবাসা এক অন্তরে থাকতে পারে।
\item \textbf{অভিশাপ নয় — সংশোধনের আশা:} মুসলিমের দায়িত্ব পাপীর জন্য দোয়া ও সংশোধনের পথ; অভিশাপ শয়তানের কাজে সাহায্য করে — পাপীকে ছেড়ে দেওয়া নয়।
\item \textbf{হদ আদায় মানে ত্যাগ নয়:} নবীজি (সা.) দণ্ড দিয়েছেন — কিন্তু তাকে ভালোবাসাও জানতেন; পাপের শাস্তি ও পাপীর প্রতি স্নেহ দুই-ই থাকতে পারে।
\item \textbf{তাওবার যাত্রা:} এই সাহাবির লড়াই দেখায় — কেউ এক রাতে তাওবাকারী হয় না; বারবার পড়ে আবার উঠে দাঁড়ানোই মুমিনের লড়াই (মুসলিম ২৭৫৯ — বারবার পাপ-তাওবা কবুল)।
\end{enumerate}
\medskip\hrule\medskip

\subsection*{৬. রেফারেন্স}

\begin{itemize}
\item সহীহ বুখারি ৬৭৮০ (কিতাবুল হুদুদ)
\item সহীহ মুসলিম ২৭৬৮
\item কুরআন: আল-ইমরান ৩:১৩৫-১৩৬
\item ব্যাখ্যা: ফাতহুল বারি — মদ্যপায়ীকে অভিশাপের আলোচনা
\end{itemize}

\end{document}
